\documentclass[a4paper,10pt]{report}\input{../0.Préambule/Préambule_cours_prof}\begin{document}

\newpage
\section*{Probabilités}
\addcontentsline{toc}{section}{Probabilités}

\begin{exer1}
Pour chaque question, entoure la (ou les) bonnes réponses.

\begin{enumerate}[font=\color{exer1}]
\item Un jeu de $32$ cartes comporte quatre rois. On tire une carte au hasard. Quelle est la probabilité d'obtenir un roi ?

\vspace{-3mm}
\begin{enumerate}[font=\color{exer1}]
\renewcommand{\arraystretch}{1}
\begin{tabularx}{\linewidth}{*{4}{>{\arraybackslash}X}}
\item $\dfrac{1}{8}$&
\item $\dfrac{1}{32}$&
\item $\dfrac{4}{32}$&
\item $\dfrac{4}{28}$\\
\end{tabularx}
\end{enumerate}
\item Une urne contient deux boules blanches, cinq boules vertes et trois boules bleues. Quelle est la probabilités d'obtenir une boule blanche ?

\vspace{-3mm}
\begin{enumerate}[font=\color{exer1}]
\renewcommand{\arraystretch}{1}
\begin{tabularx}{\linewidth}{*{4}{>{\arraybackslash}X}}
\item $\dfrac{2}{8}$&
\item $\dfrac{2}{10}$&
\item $\dfrac{8}{2}$&
\item $\dfrac{1}{5}$\\
\end{tabularx}
\end{enumerate}
\item Une urne contient cinq boules bleues numérotées de $1$ à $5$, trois boules blanches numérotées de $1$ à $3$ et deux boules noires numérotées de $1$ à $2$. Quelles est la probabilité d'obtenir un $3$ ?

\vspace{-3mm}
\begin{enumerate}[font=\color{exer1}]
\renewcommand{\arraystretch}{1}
\begin{tabularx}{\linewidth}{*{4}{>{\arraybackslash}X}}
\item $\dfrac{3}{10}$&
\item $\dfrac{2}{10}$&
\item $\dfrac{2}{5}$&
\item $\dfrac{1}{5}$\\
\end{tabularx}
\end{enumerate}
\end{enumerate}
\end{exer1}

\begin{exer1}
On considère l'urne ci-contre :

\noindent \begin{minipage}{11.5cm}
\begin{enumerate}[font=\color{exer1}]
\item Si on s'intéresse à la couleur de la boule, quelles sont les issues possibles ?
\item Si on s'intéresse au numéro écrit, quelles sont les issues possibles ?
\item Cite un événement impossible.
\item Quelle est la probabilité de l'événement A :\og Obtenir une boule rouge \fg{} 
\item Écris à l'aide d'une phrase non négative, l'événement contraire de l'événement A.
\item Quelle est la probabilité de cet événement ?
\end{enumerate}
\end{minipage}
\hspace{5mm}
\begin{minipage}{5.5cm}
\includegraphics[scale=0.2]{../40.Image/17. Exo-Probabilité 1.png}
\end{minipage}
\end{exer1}

\begin{exer1}
On considère la roue de loterie ci-contre :

\noindent \begin{minipage}{11.5cm}
\begin{enumerate}[font=\color{exer1}]
\item Si on s'intéresse aux couleurs de chaque secteur de la roue, cite les issues possibles.
\item Si on s'intéresse aux nombres de chaque secteur, cite les issues possibles.
\item Cite un événement certain de se réaliser.
\item Cite un événement impossible.
\item Quelle est la probabilité d'obtenir un nombre supérieur ou égal à $100$.
\item Écris à l'aide d'une phrase non négative l'événement contraire de l'événement \og obtenir un nombre supérieur ou égal à $100$ \fg{} 
\item Quelle est la probabilité de cet événement ?
\end{enumerate}
\end{minipage}
\begin{minipage}{5cm}
\includegraphics[scale=0.3]{../40.Image/17. Exo-Probabilité 2.png}
\end{minipage}
\end{exer1}

\vspace{-3mm}
\begin{exer1}
Précise si les affirmations suivantes sont vraies ou fausses.

\begin{enumerate}[font=\color{exer1}]
\item La probabilité de gagner à un jeu peut être égale à $1,2$.
\item Sachant que la probabilité de gagner à un jeu est égale à $0,6$, la probabilité de perdre est égale à $1,6$.
\item Sachant que la probabilité de perdre à un jeu est de $0,3$, la probabilité de gagner est égale à $0,7$.
\item Si la probabilité d'un événement est égale à $0,25$, la probabilité de l'événement contraire est égale à $0,85$.
\item La somme des probabilité d'un événement et de son événement contraire est égale à $1$.
\item La probabilité d'un événement est compris entre $0$ et $1$.
\end{enumerate}
\end{exer1}

\begin{exer2}
Il y a dans une urne douze boules indiscernables au toucher, numérotées de $1$ à $12$. On veut tirer une boule au hasard.

\begin{enumerate}[font=\color{exer2}]
\item Est-il plus probable d'obtenir un numéro pair ou bien un multiple de $3$ ?
\item Quelle est la probabilité d'obtenir un numéro inférieur à $20$ ?
\item On enlève de l'urne toutes les boules dont le numéro est un diviseur de $6$. On veut à nouveau tirer une boule au hasard. Explique pourquoi la probabilité d'obtenir un numéro qui soit un nombre premier est alors $0,375$.
\end{enumerate}
\end{exer2}

\begin{exer2}
Un dé cubique a $6$ faces peintes : une en bleu, une en rouge, une en jaune, une en vert et deux en noir. Calcule les probabilités des événement : \og Obtenir une face ...

\vspace{-2mm}
\begin{enumerate}[font=\color{exer2}]
\renewcommand{\arraystretch}{1}
\begin{tabularx}{\linewidth}{*{5}{>{\arraybackslash}X}}
\item ... jaune \fg{} &
\item ... bleue \fg{} &
\item ... rouge \fg{} &
\item ... verte \fg{} &
\item ... noire \fg{} \\
\end{tabularx}
\end{enumerate}
\end{exer2}

\begin{exer3}
Sam préfère les bonbons bleus. Dans son paquet de $500$ bonbons, $150$ sont bleus, les autres sont rouges, jaunes ou verts.

\begin{enumerate}[font=\color{exer3}]
\item Quelle est la probabilité qu'il pioche au hasard un bonbon bleu dans son paquet ?
\item $20 \%$ des bonbons de ce paquets sont rouges. Combien de bonbons rouges y-t-il ?
\item Sachant que $130$ bonbons sont verts dans ce paquet, Sam a-t-il plus de chance de piocher au hasard un bonbon vert ou un bonbon jaune ?
\end{enumerate}
\end{exer3}

\begin{exer3}

\vspace{-5mm}
\noindent \begin{minipage}{8.5cm}
Louise a téléchargé une liste de lecture sur on lecteur MP4.

\medskip
\noindent Louise décide d'utiliser la fonction \og aléatoire \fg{} de son MP4. Cette fonction choisit au hasard une chanson parmi celles qui sont présente dans la liste de lecture. 

\medskip
\noindent Chaque chanson a la même probabilité d'être écoutée.

\begin{enumerate}[font=\color{exer3}]
\item Détermine la probabilité que Louise écoute une chanson de Maen.
\vspace{5mm}
\item Elle répète $20$ fois l'utilisation de cette fonction et note à chaque fois le nom de l'interprète qu'elle a écouté. Les résultats qu'elle obtient sont notés dans le graphique ci contre.

\medskip
Quelle est la proportion de chansons de Maen écoutées ? Compare avec la question \textcolor{exer3}{\textbf{1.}}
\end{enumerate}
\end{minipage}
\hspace{5mm}
\begin{minipage}{7.5cm}
\begin{tabularx}{\linewidth}{|*{3}{>{\centering \arraybackslash}X|}}
\hline
\rowcolor{exer3!20} Titre de la chanson & Nom de l'interprète & Durée de la chanson \\\hline
Mamatéou & Tiamaté & $232$s \\\hline
La différence & Tiamaté & $211$s \\\hline
Amazing & Tiamaté & $214$s \\\hline
Tes racines & Tiamaté & $175$s \\\hline
YoungBov & Hudad & $336$s \\\hline
La ficelle & Maen & $191$s \\\hline
Fou fou fou & Maen & $184$s \\\hline
Nina & Maen & $217$s \\\hline
\end{tabularx}

\begin{center}
\begin{tikzpicture}[line cap=round,line join=round,>=triangle 45,x=1cm,y=1cm]
\begin{axis}[x=1.7cm,y=0.2cm,axis lines=middle,ymajorgrids=true,xmajorgrids=true,
xmin=0,xmax=3.5,ymin=0,ymax=17.5,
xtick={1,2,3},
xticklabels={Tiamaté, Hudad, Maen},
ytick={-0,2,...,16},
ylabel=Nb de chansosn écoutées]
\clip(0,0) rectangle (3.5,17.5);
\draw[line width=0.8pt,color=exer3,fill=exer3,fill opacity=0.3] (0.75,0.) rectangle (1.25,15);
\draw[line width=0.8pt,color=exer3,fill=exer3,fill opacity=0.3] (1.75,0.) rectangle (2.25,4);
\draw[line width=0.8pt,color=exer3,fill=exer3,fill opacity=0.3] (2.75,0.) rectangle (3.25,6);
\end{axis}
\end{tikzpicture}
\end{center}
\end{minipage}
\end{exer3}

\begin{exer3}
Un sac opaque contient $50$ billes bleues, $45$ rouges, $45$ vertes et $60$ jaunes. les billes sont indiscernables au toucher. On tire une bille au hasard dans ce sac.

\begin{enumerate}[font=\color{exer3}]
\item Donne les issues possibles de cette expérience aléatoire.
\item Cite un événement impossible
\item Quelle est la probabilité que la bille soit jaune ?
\item Calcule la probabilité de l'événement contraire à l'événement \og La bille est jaune \fg{} 
\item Calcule la probabilité de l'événement \og La bille est bleue ou verte \fg{}
\end{enumerate}
\end{exer3}

\begin{exer3}

\vspace{-5mm}
\noindent \begin{minipage}{10.5cm}
Une cible de fléchettes est composées de trois secteurs circulaires de différentes couleurs. On suppose que, pour chaque lancer, la fléchette atteint l'une des trois zones représentées ci-dessous. Les trois cercles ont le même centre. Le cercle vert a un rayon égal à $2$cm.

\begin{enumerate}[font=\color{exer3}]
\item Calcule la probabilité d'atteindre la zone verte.
\item Calcule la probabilité d'atteindre la zone rouge.
\end{enumerate}
\end{minipage}
\hspace{5mm}
\begin{minipage}{7.5cm}
\includegraphics[scale=0.3]{../40.Image/17. Exo-Probabilité 3.png}
\end{minipage}
\end{exer3}
\end{document}